\section{Conclusión}
\subsection{Primer experimento}
En el primer experimento se logró verificar que el error absoluto máximo obtenido con el multímetro digital es coherente con la clase del instrumento, lo que indica que las mediciones realizadas son confiables dentro de los límites establecidos por su clase. Esto se evidenció al calcular la clase utilizando la fórmula correspondiente y comparar el resultado con la clase indicada en el dispositivo, confirmando que estamos dentro de lo esperado.
\subsection{Tercer experimento}
En este tercer experimento se demostro que aunque no se pueda acceder a una malla conductora o tierra, el cambio que genera en la medición de la resistencia de contacto entre los electrodos no va a significar que se obtenga una conclusión errónea, ya que por ejemplo en el desarrollo de las dos mediciones una de mas del doble que la otra, en ambas se puede concliur que son buenos valores de resistencia de puesta a tierra, si no fuese así, habría que colocar jabalinas en el sistema que actuasen como resistencia en paralelo para mejorar la resistencia de puesta a tierra. 